\chapter{Introduction}

\subsection*{Cadre du stage}

Ce stage s'inscrit dans la continuité du projet eSWARM et poursuit les travaux d'Augustin Bonnel \cite{bonnel_modeling_2024} et de Sylvain Boegeat, deux anciens étudiants de l'INSA. L'objectif de ce projet est de développer un système de drones quadricoptères modulaire et reconfigurable, capable de s'assembler en vol afin d'accomplir des missions qu'un drone seul ne pourrait réaliser, comme par exemple le transport de charges lourdes.



\subsection*{Définition}

Le projet eSWARM est basé sur des essaims de robots dits reconfigurables et modulaires. Une définition de ce type de système est donnée dans \cite{seo_modular_2019} : "Les systèmes de robots modulaires reconfigurables (MRR pour \textit{Modular Reconfigurable Robot}) sont constitués de nombreux modules (ou unités) répétés qui peuvent être réorganisés dans différentes configurations en fonction de la tâche que le robot doit accomplir." Une définition similaire est fournie dans \cite{brooks_overview_2021}.

Il est essentiel d'expliquer les deux termes : \textit{modulaire} et \textit{reconfigurable}. Le premier désigne une architecture composée de modules pouvant s'assembler de diverses manières. Le second caractérise la capacité du système à modifier sa structure au cours de son fonctionnement.

Il existe des systèmes modulaires reconfigurables terrestres, tels que le robot SABER \cite{noauthor_saber_nodate} ou encore PolyBot \cite{yim_modular_nodate}. Ce dernier se compose de modules, comme des segments articulés et des nœuds de connexion, pouvant être assemblés pour former différentes structures. Il peut changer de forme automatiquement, par exemple en passant d'une configuration en forme de serpent pour explorer des terrains complexes à une structure plus rigide pour manipuler des objets. 

En ce qui concerne les systèmes modulaires reconfigurables aériens, composés d'UAV (\textit{Unmanned Aerial Vehicles}), de nombreux documents y font référence. Par exemple, \cite{zhang_design_2024}, \cite{saldana_modquad_2018}, et \cite{gabrich_modquad-dof_2020} présentent des projets similaires à eSWARM, avec des drones quadricoptères capables de s'attacher en vol pour former une structure unique et pilotable. Des détails supplémentaires sur ces exemples seront présentés dans la partie \ref{essaim_de_drones} du présent rapport.

